电子初态的有限纵向带宽则通过输入密度矩阵进入。令
\begin{equation*}
 \rho_\ee^{\rm in}
 =\int\!\dd p\,\dd p'\,
 \rho_e(p,p')|p\rangle\langle p'|,
\end{equation*}
并用总散射算符 $S$ 定义电子约化输出
\begin{equation}
\boxed{
 \rho_\ee^{\rm out}
 =\operatorname{Tr}_{\rm ph}
 \left[S(\rho_\ee^{\rm in}\otimes\rho_{\rm ph})S^\dagger\right].}
\label{eq:finite_envelope_reduced_output_dp47}
\end{equation}
若探测器只分辨相对于局部中心动量的边带阶数，有限能量分辨函数 $R_E(\Delta E)$ 进一步给出实验概率
\begin{equation}
\boxed{
 P_j^{\rm meas}
 =\int\!\dd E\,R_E(E-E_j)
 \langle E|\rho_\ee^{\rm out}|E\rangle,
 \qquad
 \int\!\dd E\,R_E(E)=1 .}
\label{eq:pinem_detector_convolution_dp47}
\end{equation}
卷积映射把理论边带概率与谱仪像素计数区分开；后续统计反演应使用\eqnrefs{eq:pinem_detector_convolution_dp47} 的卷积核，而不是直接把理想 $P_\ell$ 与实验直方图比较。

有限输入能散 $\sigma_E$ 是否会洗掉相邻边带取决于两个不同条件。分辨边带需要
\begin{equation}
 \epsilon_{\rm res}\equiv\frac{\sigma_E}{\hbar\omega}\ll1,
\end{equation}
而保持输入相干的相位条件由输入密度矩阵非对角宽度 $\sigma_{E,{\rm coh}}$ 决定。二者一般不相等；能谱窄并不自动意味着长相干时间。
